\section{Introduction}
\label{sec:introduction}

The interpretation of genomic variants remains a central bottleneck in clinical
genetics and precision medicine. Despite advances in sequencing technology,
the majority of rare variants detected in human genomes are classified as
variants of uncertain significance (VUS), limiting their clinical
actionability~\cite{ClinVar_2018,ACMG_2015}. Computational variant effect
prediction (VEP) tools provide evidence for the ACMG/AMP PP3 and BP4 criteria,
but each existing tool addresses only a subset of the variant landscape.
AlphaMissense~\cite{AlphaMissense_2023} and REVEL~\cite{REVEL_2016} are
restricted to missense variants; SpliceAI~\cite{SpliceAI_2019} predicts only
splice-altering effects; non-coding scores such as ncER~\cite{ncER_2019} and LINSIGHT~\cite{LINSIGHT_2017}
provide regional constraint estimates but lack variant-level resolution;
and integrative scores such as CADD~\cite{CADD_2024},
though broadly applicable, incorporate ClinVar annotations in their training,
creating potential circularity when used for ClinVar variant
reclassification~\cite{gnomAD_2020}. No single existing tool provides unified,
training-data-free predictions across coding, non-coding, and indel variants.

Genomic foundation models offer a fundamentally different approach. Trained on
large collections of DNA sequences in a self-supervised manner, these models
learn the statistical grammar of genomes without exposure to phenotypic labels.
The Evo family of models---Evo~\cite{Evo1_2024} and
Evo2~\cite{Evo2_2026}---represents the current state of the art in genomic
language modelling, with Evo2's 7-billion-parameter architecture trained on 9.3
trillion nucleotides from 128,000 genomes spanning all domains of life. Evo2
supports context windows up to one million base pairs, far exceeding the Nucleotide
Transformer~\cite{NT_2024} ($\sim$6\,kb effective context) and
HyenaDNA~\cite{HyenaDNA_2023} (450\,kb at reduced resolution). The zero-shot
variant effect prediction paradigm---comparing reference and alternate sequence
log-likelihoods without fine-tuning---enables variant scoring that is intrinsically
free of training-set circularity and agnostic to variant type, including
non-coding and indel variants that remain inaccessible to most supervised tools.

Spaceflight presents a unique context for variant interpretation. Astronauts are
exposed to galactic cosmic radiation (GCR) and solar particle events that
generate characteristic mutational signatures, including C$>$A transversions
from 8-oxoguanine lesions and complex double-strand break (DSB)
repair products~\cite{daSilveira_2020}. The NASA Twins Study revealed persistent
gene expression changes in DNA repair, telomere maintenance, and immune
regulation pathways after long-duration spaceflight~\cite{Twins_2019}, and
Rutter et al.~\cite{Rutter_2024} recently identified protective and risk alleles
in astronaut genomes across genes involved in radiation response. However, the
functional consequences of the vast majority of variants in these genes---particularly
non-coding and structural variants---remain unknown. A comprehensive variant
effect atlas for spaceflight radiation-response genes would support both
mission-planning risk assessment and fundamental understanding of
radiation-induced mutagenesis.

Here we apply Evo2 to generate zero-shot variant effect predictions for 215,001
variants (168,409 SNVs and 46,592 indels) across ten genes selected for their
relevance to spaceflight radiation biology: four well-characterised control genes
(BRCA1, TP53, CHEK2, DNMT3A) with extensive experimental validation data, and
six spaceflight-associated genes (TERT, ATM, NFE2L2, CLOCK, MSTN, RAD51). We
validate predictions through a three-layer framework: (1) deep mutational
scanning (DMS) calibration against four independent functional datasets, (2)
ClinVar clinical benchmarking with multi-tool comparison and matched-variant
analysis, and (3) non-coding validation using TERT promoter MPRA data and ENCODE
regulatory element annotations.

We demonstrate that Evo2 achieves clinical-grade ClinVar AUROCs (0.91--1.00),
is the highest-performing tool for two of six benchmarked genes, provides
orthogonal signal that improves ensemble classification, and---critically---is
the only tool with significant positive correlation with non-coding MPRA data
in the TERT promoter. The resulting variant effect atlas, encompassing coding
and non-coding regions across all ten genes, is provided as a public resource.
